\documentclass[11pt]{article}

%==============================================================================
%  LDS -- Local Dimension and Scale: Design and Test Plan
%  Phase spec for the geosmooth program. Written to be implementable without
%  gaps and auditable. Plain language; every term defined on first use.
%==============================================================================

\usepackage[margin=1in]{geometry}
\usepackage{amsmath,amssymb}
\usepackage{xcolor}
\usepackage{booktabs}
\usepackage{array}
\usepackage{enumitem}
\usepackage{microtype}
\microtypesetup{expansion=false}
\usepackage{caption}
\usepackage[most]{tcolorbox}
\usepackage{datetime2}
\DTMsetup{showzone=false}
\usepackage[hidelinks]{hyperref}
\widowpenalty=10000
\clubpenalty=10000
\raggedbottom

\definecolor{ink}{HTML}{24303D}
\definecolor{accent}{HTML}{2563EB}
\definecolor{accent3}{HTML}{15803D}
\definecolor{accent4}{HTML}{B45309}
\setlist{topsep=3pt,itemsep=2pt,parsep=1pt}

\newtcolorbox{gate}[1]{
  enhanced, colback=accent3!5, colframe=accent3!55!black, boxrule=0.4pt, arc=2pt,
  left=7pt, right=7pt, top=4pt, bottom=4pt, title={#1},
  fonttitle=\bfseries\sffamily\small, coltitle=accent3!42!black,
  attach title to upper={\quad}}
\newtcolorbox{decide}{
  enhanced, colback=accent4!6, colframe=accent4!55!black, boxrule=0.4pt, arc=2pt,
  left=7pt, right=7pt, top=4pt, bottom=4pt, title=Decision needed from the orchestrator,
  fonttitle=\bfseries\sffamily\small, coltitle=accent4!50!black,
  attach title to upper={\quad}}
\newtcolorbox{plain}{
  enhanced, colback=accent!4, colframe=accent!55!black, boxrule=0.4pt, arc=2pt,
  left=7pt, right=7pt, top=4pt, bottom=4pt, title=In plain words,
  fonttitle=\bfseries\sffamily\small, coltitle=accent!50!black,
  attach title to upper={\quad}}

\newcommand{\field}[1]{\smallskip\noindent\textbf{\textsf{#1.}}\ }
\newcommand{\code}[1]{\texttt{#1}}

\title{\textbf{LDS --- Local Dimension and Scale}\\[3pt]
\large Design and Test Plan (implementation-grade)\\[5pt]
\normalsize Make the neighborhood size local, and choose it together with the local
dimension. The raw estimator; spatial smoothing of the two fields is a later phase.}
\author{\small geosmooth / LPS program\\[2pt]
\footnotesize \nolinkurl{project_briefs/lps_lds_design_and_test_plan_2026-06-13.tex}}
\date{\DTMnow\ \textnormal{\footnotesize ET}}

\begin{document}
\maketitle

\begin{abstract}\noindent
Today the method uses one neighborhood size for the whole dataset and, separately,
guesses how many directions the data spreads in (the ``local dimension'') at each
point using that one fixed size. This phase, \textbf{LDS} (Local Dimension and
Scale), makes the neighborhood size \emph{local} --- different at different points
--- and chooses it \emph{together} with the local dimension, because the two affect
each other: a neighborhood that is too wide sees structure that is not local, too
narrow sees only noise, and the dimension you read off changes with the width. The
output is, at every point, a neighborhood size and a dimension. These are
\emph{raw}: each point decides on its own, with no smoothing across points. Smoothing
the two fields is the job of the next phases (LDR, LDSR); LDS produces the fields they
will smooth. This document is written so an implementer can build LDS without filling
any gaps, and an auditor can check it. Every technical term is defined on first use.
\end{abstract}

{\small\tableofcontents}

%==============================================================================
\section{Plain-language glossary (read once)}\label{sec:glossary}
%==============================================================================

\begin{description}[leftmargin=0em,style=nextline]
\item[Anchor.] A point at which we build one small local model. Today the anchors
are the data points themselves.
\item[Neighborhood (support).] The set of data points closest to an anchor that the
local model uses. Its \emph{size} is the count $K$ (how many points) and equivalently
a \emph{radius} $h$ (the distance to the farthest of those $K$ points).
\item[Scale.] How wide the neighborhood is. ``Choosing the scale'' means choosing
$K$ (or $h$). A larger scale uses more, farther points.
\item[Local dimension.] How many independent directions the data genuinely spreads
in, close to the anchor. On a curved sheet sitting in 3-D space the local dimension
is $2$, even though the points have $3$ coordinates. We will write it $d$.
\item[Local chart.] A small local coordinate system at the anchor, built by finding
the few directions of largest spread (principal directions) of the neighborhood. The
local model is fit in these coordinates.
\item[Degree $p$.] The order of the local polynomial fit ($p=1$ is a local straight
fit/plane; $p=2$ adds curvature). Fixed per candidate.
\item[Singular values.] Numbers $\sigma_1\ge\sigma_2\ge\cdots\ge0$ measuring how much
the neighborhood spreads along each principal direction. Big then small means the
data is essentially low-dimensional.
\item[Effective sample size (ESS).] With weighted points, the ``honest'' number of
points contributing. If the weights are $w_1,\dots,w_K$, then
$\mathrm{ESS}=\left(\sum_j w_j\right)^2/\sum_j w_j^2$. It is at most $K$ and drops when
a few points dominate. More ESS $=$ less noise in the local estimate.
\item[Linear smoother.] A method whose prediction at each point is a fixed weighted
average of the responses, $\hat y = S\,y$, where the weight matrix $S$ depends only on
\emph{where} the data points are (their inputs $X$), not on the response values $y$.
Why it matters: when this holds, exact, cheap formulas exist for the uncertainty and
the ``degrees of freedom'', and the program's existing toolkit applies unchanged. This
holds \emph{conditional on a fixed configuration}; when the scale is chosen from the
data (as it already is, globally, by cross-validation), \S\ref{sec:conditional} states
exactly what survives.
\item[Degrees of freedom (df).] A single number measuring how flexible the fit is,
equal to $\mathrm{tr}\,S$ (the sum of the diagonal of $S$). Two methods are compared
fairly only at the \emph{same} df.
\item[Self-influence $S_{ii}$.] How much a point's own response affects its own
prediction. Near $1$ means the fit nearly interpolates that point.
\item[Leave-one-out (LOO) error.] The error a method makes at a point when that
point is not used to fit it --- an honest out-of-sample error. For a linear smoother
there is a cheap shortcut: $\big(y_i-\hat y_i\big)/\big(1-S_{ii}\big)$.
\item[Bias and variance.] Bias is systematic error from over-smoothing (a too-wide
neighborhood flattens real structure). Variance is random error from under-smoothing
(a too-narrow neighborhood chases noise). The best scale trades them off.
\end{description}

%==============================================================================
\section{What exists today (the starting point)}\label{sec:today}
%==============================================================================

The implementer should read these existing pieces first; LDS extends them, it does
not replace them.

\begin{itemize}[leftmargin=1.4em]
\item \textbf{One global scale, chosen by cross-validation.} The fitter tries a few
candidate neighborhood sizes (\code{support.grid}, e.g.\ $\{10,15,20\}$) crossed with
degree and kernel, scores each by cross-validated error, and picks \emph{one}
$K$ for the whole dataset (\code{.klp.cv.table} $\to$ \code{.klp.select.best.idx} $\to$
the single \code{selected\$support.size}). That one $K$ is used at every anchor
(\code{.klp.predict.local.polynomial}).
\item \textbf{A per-point dimension, at that fixed scale.} With
\code{chart.dim = "local.auto"}, the code already estimates a dimension at each anchor
(\code{.klp.resolve.prediction.chart.dim} $\to$
\code{.klp.local.auto.chart.dim.from.order}), returning a vector
\code{chart.dim.by.eval}. The estimate uses the singular values of the neighborhood
(\S\ref{sec:dimrule}).
\item \textbf{The points-per-parameter cap.} The dimension is capped so the local
polynomial is fittable: you need at least $\binom{d+p}{p}$ points to fit a degree-$p$
polynomial in $d$ dimensions, so the code requires $\binom{d+p}{p}\le K$
(\code{.local.pca.max.chart.dim.for.support}). This is the one place scale and
dimension are already tied.
\item \textbf{The uncertainty toolkit.} \code{R/lps\_uncertainty.R} can produce the
smoother matrix $S$, its diagonal $S_{ii}$, $\mathrm{df}=\mathrm{tr}\,S$, and bands,
for a fixed configuration. LDS reuses this to score scales cheaply.
\item \textbf{A per-anchor bandwidth knob already exists.} The bandwidth multiplier
(\code{bandwidth.multiplier.grid}) added in the Tier-1 work lets the bandwidth be
scaled; LDS will set it (or $K$) per anchor.
\end{itemize}

\begin{plain}
Right now: the method decides ``use the 20 nearest points everywhere,'' then at each
point guesses the dimension from those 20. LDS changes the first part: decide the
number of nearest points \emph{per point}, and decide it jointly with the dimension.
\end{plain}

%==============================================================================
\section{The one decision that shapes everything: how the scale is chosen}\label{sec:fork}
%==============================================================================

The neighborhood size has to be chosen somehow, and there are two honest criteria.
First, what is \emph{not} in question: a geometry floor --- enough points to fit the
local polynomial (the points-per-parameter cap, \S\ref{sec:dimrule}) and a
well-conditioned design (\S\ref{sec:criteria}) --- is enforced for every candidate size.
The real decision is which criterion \emph{picks among the geometry-feasible sizes}, and
the machinery for both is built behind a switch.

\field{Option B --- ``error-driven'' scale (the localization of what we do today)}
Choose the size at each point to make the local \emph{prediction error} small, using the
response (local cross-validation / GCV, or Lepski's rule, \S\ref{sec:criteria}). This is
the faithful \emph{localization of the current method}: today's single global $K$ is
already chosen by cross-validation, which reads $y$; Option~B applies that same principle
per point instead of once for the whole dataset. It is also the classical move
(Friedman's variable-span supersmoother; Fan--Gijbels variable bandwidth; Lepski), and
it is the criterion that adapts to how \emph{wiggly} the function is --- tighter
neighborhoods where it changes fast --- which is the main reason to localize the scale
at all. Cost: the choice depends on $y$, so the method is not a fixed linear smoother;
df and bands are computed the selection-aware way (\S\ref{sec:conditional}). That is the
\emph{same} regime today's global CV already lives in, only more local --- more
selections, hence a larger selection-inflation, a difference of degree not of kind
(\S\ref{sec:conditional}). One caution: \emph{naive} per-point CV is unstable (a
near-single-residual criterion), so the robust forms are the smoothed-CV supersmoother
idea or Lepski's rule --- \textbf{Mode~B--Lepski is the recommended local mode}.

\field{Option A --- ``geometry-only'' scale (the conservative floor and fallback)}
Choose the size at each point using \emph{only the positions of the data points} --- hold
a target effective sample size (\S\ref{sec:criteria}) --- never the response. Its appeal
is that the per-anchor choice ignores $y$, so the fit stays a \emph{linear smoother} given
its target and the exact toolkit applies (\S\ref{sec:conditional}). But it adapts only
\emph{geometrically} --- tightening the radius $h$ where points are dense and raising the
count only where the feasibility floor demands it (G8) --- not to the function's wiggliness,
so it forgoes the main adaptivity prize; and per \S\ref{sec:conditional} its linearity is
\emph{conditional-only} anyway
(the default target derives from the global $y$-selected pick) --- one global selection
away from B in degree, not a difference in kind. Option~A is the right \emph{floor and
fallback}: the minimum well-posed size under both modes, the mode to use when the exact
toolkit is required, and the safe landing for a degenerate anchor.

\begin{decide}
\textbf{D1 --- recommended local mode, and how the shipped default is set.} The API
default is unchanged: a single global $K$, bit-for-bit (\S\ref{sec:impl}, I1) --- this
decision is only about the recommended mode \emph{when per-point scale is enabled}.
Recommended front-runner: \textbf{Mode~B--Lepski} (local error-driven, stabilized), with
Mode~A (geometry-only) as the exact-toolkit fallback and the well-posedness floor under
both. The \emph{shipped} default is not pre-judged here --- studies S1/S2 decide it on
D-VARSCALE at \emph{matched df} (accuracy gain), with the oracle gap (S4) as a diagnostic,
in phase LDS2 (\S\ref{sec:sequence}); the unconditional GDF / coverage cost (S3) is a later
LDS3 consideration, not an early gate. If the error-driven gain is marginal, A is the safe
default; otherwise B--Lepski ships.
\end{decide}

%==============================================================================
\section{What ``linear smoother'' guarantees --- conditional vs.\ unconditional}\label{sec:conditional}
%==============================================================================

The phrase ``linear smoother'' (\S\ref{sec:glossary}) carries an exact-toolkit promise
--- df $=\mathrm{tr}\,S$, the leave-one-out shortcut, uncertainty bands --- and it is
worth saying precisely \emph{when} that promise holds, because the scale is
\emph{chosen} from the data and choosing is not a linear operation. The test plan
(\S\ref{sec:tests}) and decision D6 (\S\ref{sec:decisions}) refer back to this section.

\field{Conditional on a fixed configuration, the fit is exactly linear}
Fix the neighborhood size, degree, kernel, and the per-anchor fields $(K_i,d_i)$. Then
the prediction is $\hat y=S\,y$ with $S$ built entirely from the inputs $X$ --- the
neighborhoods, the local charts, the kernel weights, and (under \code{local.auto}) the
singular values --- none of which read $y$. So $S$, $\mathrm{tr}\,S$, the $S_{ii}$, and
the bands are \emph{exact for that configuration}. (This is a statement about the
Gaussian/weighted-least-squares fit; the logistic/binomial fit is non-linear in $y$
regardless of selection.)

\field{Selecting the scale makes the whole pipeline non-linear}
The selected size is itself a function of the responses, $\hat K=\hat K(y)$, so the
end-to-end map is $\hat y=S\big(\hat K(y)\big)\,y$: the matrix is data-chosen and the
composition is \emph{not} linear in $y$. \textbf{This is already true today}, before
LDS --- the single global $K$ is picked by cross-validation, which reads $y$. LDS does
not introduce the issue; it only changes \emph{how many} selections happen and how
local they are.

\field{The non-linearity is piecewise, so the conditional toolkit is locally exact}
The selection criterion (CV error, local GCV, Lepski) is a continuous, piecewise-smooth
function of $y$, and $\hat K$ is its arg-best over a \emph{finite} grid. So
$y\mapsto\hat K(y)$ is a step function --- locally constant in $y$ except on a
measure-zero set where two candidate sizes tie --- and $y\mapsto\hat y$ is
\emph{piecewise linear}. At the observed $y$ (away from ties) the local sensitivity
$\partial\hat y/\partial y$ is exactly $S(\hat K)$. So $S(\hat K)$ is the exact
\emph{local} influence matrix and $\mathrm{tr}\,S(\hat K)$ the exact \emph{conditional}
degrees of freedom: the toolkit computes the right \emph{conditional} object, not a
wrong one.

\field{What is lost is unconditional honesty, not local linearity}
Across repeated samples $\hat K$ fluctuates, adding a between-configuration variance
component that $\sigma^2 SS^\top$ ignores. So $\mathrm{tr}\,S(\hat K)$ \emph{understates}
the effective degrees of freedom of the adaptive procedure, and the conditional bands
are \emph{optimistically narrow}: they answer ``if we had pre-specified this scale,''
not ``what did choosing the scale from this data cost.'' The formally correct
unconditional notions are the \emph{generalized degrees of freedom}
$\mathrm{GDF}=\sum_i \partial\hat y_i/\partial y_i$ (Ye, 1998, \emph{JASA}) and the
equivalent covariance-penalty df $\sum_i \mathrm{cov}(\hat y_i,y_i)/\sigma^2$ (Efron,
2004, \emph{JASA}); both equal $\mathrm{tr}\,S$ for a fixed linear smoother and
\emph{exceed} it once selection is in the loop.

\field{How much it bites, and what this says about the two modes}
The gap grows with the number and fineness of the selections and the flatness of the
criterion (frequent near-ties $=$ unstable selection $=$ larger selection variance).
This orders the modes:
\begin{itemize}[leftmargin=1.4em]
\item \textbf{Mode A (geometry-only)} reads no $y$ in its per-anchor choice, so
\emph{given its target} it is exactly conditionally linear, and G2 verifies it. But its
default target is the ESS of the global CV pick (\S\ref{sec:criteria}, D2), so the
pipeline still inherits \emph{one} global, coarse, $y$-dependent selection --- the same
one today's method already carries. Mild, and unchanged from the status quo.
\item \textbf{Mode B (error-driven)} adds a \emph{per-anchor}, $y$-driven selection at
every point; this is where the unconditional gap can become non-trivial. It is the
price of the accuracy benefit, and it must be reported the selection-aware way (I6, S2,
S3).
\end{itemize}

\begin{plain}
Once you fix how wide each neighborhood is, the method is a clean weighted average and
every exact formula applies. But the width is chosen from the data --- already today,
globally; and at every point if we let each point choose --- and choosing isn't a
linear step. So the usual df and error bars are exactly right \emph{for the width that
was picked}, and a little too optimistic as a statement about the whole pick-then-fit
procedure. The fix is either to say plainly that the bars are conditional on the chosen
width, or to measure the extra cost of choosing (GDF / covariance penalty, or an outer
split).
\end{plain}

\noindent\textbf{Practical resolutions} (the menu D6 picks from): (i) report
conditionally and label it; (ii) split selection from inference --- choose the scale on
one fold, fit and infer on another (the nested-CV logic); (iii) estimate the inflated
df directly by the Ye perturbation / Efron covariance penalty (jitter $y$ by small
i.i.d.\ noise, refit \emph{including reselection of the scale}, read off
$\sum_i\mathrm{cov}(\hat y_i,y_i)/\sigma^2$); (iv) bootstrap the whole pipeline
including reselection. Mode~A defaults to (i); Mode~B should additionally report (iii)
so the cost is visible (S3).

%==============================================================================
\section{The estimator, step by step}\label{sec:algorithm}
%==============================================================================

This is the complete per-anchor procedure. Inputs: training inputs $X$ ($n$ points,
$D$ ambient coordinates), the anchor center $c$, a candidate scale grid
$\mathcal{K}$ (\S\ref{sec:grid}), the degree $p$, the kernel, the design basis, the
scale mode (A or B), and the mode's target/criterion. Output: a chosen size $K^\star$
and dimension $d^\star$ at this anchor, plus diagnostics.

\begin{enumerate}[leftmargin=1.5em]
\item \textbf{Order by distance.} Compute the distance from $c$ to every training
point and sort. (One sort serves all candidate sizes.)
\item \textbf{For each candidate size $K\in\mathcal{K}$ with $K\le n$:}
  \begin{enumerate}[leftmargin=1.4em,label=(\alph*)]
  \item \textbf{Take the neighborhood.} The $K$ nearest points; set the radius
  $h=$ distance to the $K$-th nearest; weights $w_j=\text{kernel}(\text{dist}_j/h)$.
  \item \textbf{Read the spread.} Center the $K$ points at $c$; take the singular
  values $\sigma_1\ge\sigma_2\ge\cdots$ of the centered block.
  \item \textbf{Estimate the dimension $d(K)$} by the rule of \S\ref{sec:dimrule}
  (which already enforces the points-per-parameter cap for this $K$ and $p$).
  \item \textbf{Score the size} by the chosen mode:
    \begin{itemize}[leftmargin=1.2em]
    \item Mode A: $\text{score}(K)=\big|\mathrm{ESS}(K)-\mathrm{ESS}_{\text{target}}\big|$,
    among sizes that pass the well-posedness filter (\S\ref{sec:criteria}).
    \item Mode B (local GCV, \emph{diagnostic only}): $\text{score}(K)=$ the local
    leave-one-out error at the anchor for the fit at $(K,d(K))$ (\S\ref{sec:criteria});
    the scale-field decision rule is Lepski, next.
    \item Mode B (Lepski): handled by the separate widen-until-it-moves rule
    (\S\ref{sec:criteria}); it does not fit the simple ``score then argmin'' shape.
    \end{itemize}
  \end{enumerate}
\item \textbf{Pick.} $K^\star=$ the chosen size: \textbf{Mode~A} --- the smallest ESS-gap
score; \textbf{Mode~B} --- Lepski's largest accepted size. (Local GCV is computed only as a
diagnostic, never as the decision.) $d^\star=d(K^\star)$.
\item \textbf{Guard.} If no size passes the filter (e.g.\ a degenerate anchor),
fall back to the global size and the summary dimension, and record the fallback in
the diagnostics (never silently).
\end{enumerate}

\begin{plain}
For each point: sort the others by closeness; try a handful of neighborhood sizes;
for each size read off how many directions the data spreads in; keep the size that
best meets the rule (geometry target, or the Lepski stability rule); report that size and
that dimension.
\end{plain}

%==============================================================================
\section{The scale-selection criteria, written out}\label{sec:criteria}
%==============================================================================

\subsection{Mode A --- geometry-only (effective sample size target)}
At each anchor and size $K$, compute $\mathrm{ESS}(K)=\big(\sum_j w_j\big)^2/\sum_j w_j^2$
from the kernel weights at scale $K$. Choose the smallest $K\in\mathcal{K}$ with
$\mathrm{ESS}(K)\ge \mathrm{ESS}_{\text{target}}$; if none reaches the target, take
the largest available $K$. Apply two hard filters first: (i) points-per-parameter on the
\emph{effective} sample size, $\mathrm{ESS}(K)\ge \rho\binom{d(K)+p}{p}$ --- \emph{not} raw
$K$: with tricube/Gaussian weights $\mathrm{ESS}(K)$ can fall well below $K$, so a raw-$K$
feasibility test over-promises (a $K=30$ neighborhood with heavy down-weighting may carry
$\mathrm{ESS}=12$, too few to fit a degree-2 surface in 3-D stably), with safety factor
$\rho\ge1$ (default $1.5$; \S\ref{sec:decisions}, D7) --- this refines the existing raw-$K$
cap; and (ii) well-conditioned design --- the local polynomial design matrix in the chart
has condition number $\le$ \code{ridge.condition.max}; if not, that $K$ is skipped.

\begin{plain}
``Use just enough nearby points to get a target amount of effective data, and at
least enough to fit the local shape stably.'' Where points are sparse the radius
grows automatically; where dense it stays tight. The noise level of the estimate is
roughly constant across the data.
\end{plain}

\noindent Default $\mathrm{ESS}_{\text{target}}$: the ESS of the global
cross-validation-selected $K$ (so on uniform data LDS reproduces today's behavior),
or a fixed value the orchestrator pins (\S\ref{sec:decisions}, D2).

\subsection{Mode B --- error-driven, local GCV (a diagnostic, not the first decision rule)}
Fit the local model at $(K,d(K))$. The fit is a weighted average
$\hat m(c)=\sum_j s_j\,y_j$; the self-influence of the anchor's own point is
$s_{\text{self}}$ (the entry of $S$ for the anchor on itself). The local
leave-one-out error proxy is $e(K)=\big(y_{\text{self}}-\hat m(c)\big)/\big(1-s_{\text{self}}\big)$;
take $\text{score}(K)=e(K)^2$ (optionally averaged over the anchor and its few nearest
neighbors for stability). Choose the $K$ minimizing it.
\emph{Status:} per \S\ref{sec:conditional} and the instability of this
near-single-residual criterion, local GCV is implemented as a \emph{reported diagnostic}
and research comparator, \emph{not} the scale-field decision rule. The decision rule is
Mode~B--Lepski below; the honest-LOO caveat is the reason, and the sequencing
(\S\ref{sec:sequence}) builds Lepski first.

\begin{gate}{Honest caveat (do not skip --- the auditor checks this)}
For a $K$-nearest-neighborhood, \emph{removing} the anchor's own point changes which
points are the $K$ nearest (the $K$-th neighbor shifts), so the cheap shortcut above
is slightly optimistic for \emph{scale} selection. Mitigate one of three ways, and
state which: (i) recompute the neighborhood without the anchor (true leave-one-out)
for the scored points only; (ii) leave a small block out within the neighborhood; or
(iii) prefer Mode B--Lepski below, which compares sizes directly and avoids the
shortcut entirely.
\end{gate}

\subsection{Mode B --- error-driven, Lepski (recommended for the scale field)}
For increasing $K$, compute the local estimate $\hat m_K(c)$ and its standard error
$\mathrm{SE}_K=\hat\sigma\,\lVert s_{\cdot}\rVert$ (from the toolkit).
\emph{Pinned project variant} (the auditor checks this exact rule, not a re-interpretation):
process the grid $K_1<K_2<\cdots$ from the smallest; call $K_j$ \emph{admissible} if it
agrees with \emph{every} smaller grid size,
$\big|\hat m_{K_j}-\hat m_{K_i}\big|\le \tau\,(\mathrm{SE}_{K_j}+\mathrm{SE}_{K_i})$ for all
$i<j$; \emph{stop at the first inadmissible size} and set $K^\star$ to the largest size
reached before it (the monotone Lepski stopping rule). If even $K_2$ is inadmissible,
$K^\star=K_1$. The tolerance $\tau$ is fixed (default $\tau=2$; \S\ref{sec:decisions}, D8).

\smallskip\noindent The noise level $\hat\sigma$ in $\mathrm{SE}_K$ must be \emph{pinned}:
Lepski's choice is sensitive to it and to $\tau$. Default --- a \emph{global, robust}
$\hat\sigma$ from a pilot global-$K$ fit (a MAD- or difference-based residual-scale
estimate, homoscedastic assumption); in the synthetic gates use the \emph{known} $\sigma$
so the test isolates scale selection from $\hat\sigma$-estimation error. Local /
heteroscedastic $\hat\sigma$ is a later extension (\S\ref{sec:decisions}, D8). The
$S$-studies report sensitivity to $(\hat\sigma,\tau)$.

\begin{plain}
``Keep widening the neighborhood until widening it would change the answer by more
than the noise allows; stop just before.'' This balances over- and under-smoothing at
each point on its own, and it naturally re-reads the dimension at every width.
\end{plain}

%==============================================================================
\section{The dimension rule at a given scale (reuse, stated exactly)}\label{sec:dimrule}
%==============================================================================

This is the existing rule (\code{.local.pca.auto.chart.dim.from.singular.values}); LDS
calls it once per candidate size, with \emph{one refinement} --- the feasibility cap
(step~2) uses the \emph{effective} sample size $\mathrm{ESS}(K)$ rather than raw $K$
(\S\ref{sec:criteria}, D7). \textbf{This refinement is gated to the LDS local-scale path
only}: the existing global-$K$ \code{chart.dim = "local.auto"} path keeps the raw-$K$ cap
unchanged, so the bit-for-bit default preservation (I1) is untouched. Given positive
singular values $\sigma_1\ge\cdots$, neighborhood size $K$, degree $p$, ambient $D$:

\begin{enumerate}[leftmargin=1.5em]
\item Energy $\epsilon_j=\sigma_j^2$; total $T=\sum_j \epsilon_j$.
\item Points-per-parameter cap (ESS-based): the largest $d$ with
$\binom{d+p}{p}\le \mathrm{ESS}(K)/\rho$ (the effective-sample-size feasibility of
\S\ref{sec:criteria}, safety factor $\rho$); call it $d_{\text{cap}}$. The existing code
uses $K-1$; LDS substitutes $\mathrm{ESS}(K)/\rho$.
\item Maximum allowed dimension: $d_{\max}=\min\big(D,\ d_{\text{cap}},\ \#\{\sigma_j\},\ K-1\big)$.
\item Variance dimension $d_{\text{var}}=$ the smallest $d$ with
$\big(\sum_{j\le d}\epsilon_j\big)/T\ge 0.95$, clamped to $[1,d_{\max}]$.
\item Eigengap dimension: among $j=1,\dots,d_{\max}-1$ find the largest ratio
$\sigma_j/\sigma_{j+1}$; if that largest ratio is $\ge 4$, $d_{\text{gap}}=$ its index,
else none.
\item Selected $d=d_{\text{var}}$ if there is no eigengap, else
$\min(d_{\text{var}},d_{\text{gap}})$; clamp to $[1,D]$.
\end{enumerate}

\begin{plain}
``Keep enough directions to explain 95\% of the spread, but if there is a sharp drop
in spread (a factor of $4$) cut there instead; never keep more directions than the
neighborhood can support a fit for.''
\end{plain}

%==============================================================================
\section{The candidate scale grid}\label{sec:grid}
%==============================================================================

A geometric grid of neighborhood sizes, wide enough to bracket the per-anchor optimum
and fine enough to resolve it. Default $\mathcal{K}=\{12,18,27,40,60,90,135\}$
(roughly $\times1.5$ steps), clipped to $[\,d_{\min}\text{-feasible},\,n-1\,]$. The
lower end must always satisfy points-per-parameter for $d=1$ and the degree in use.
The grid is a parameter; the orchestrator may pin it (D3).

%==============================================================================
\section{Outputs and diagnostics}\label{sec:outputs}
%==============================================================================

LDS returns, alongside the prediction, two per-anchor vectors and a diagnostics
table:
\begin{itemize}[leftmargin=1.4em]
\item \code{support.size.by.eval} --- the chosen $K^\star_i$ at each anchor (new).
\item \code{chart.dim.by.eval} --- the chosen $d^\star_i$ at each anchor (already
exists; now produced jointly with $K$).
\item per-anchor diagnostics: the scored sizes, the realized ESS, the dimension rule
that fired (variance vs eigengap), $s_{\text{self}}$, any fallback flag, the mode used.
\end{itemize}
These two raw fields are the input to the later smoothing phases (LDR, LDSR).

%==============================================================================
\section{Implementation plan (API, functions, defaults)}\label{sec:impl}
%==============================================================================

\begin{enumerate}[leftmargin=1.5em]
\item[I1.] \textbf{New mode value.} Allow the neighborhood size to be a per-anchor
field: a new accepted value \code{support.size = "local.auto"} (parallel to
\code{chart.dim = "local.auto"}), plus arguments \code{scale.mode = c("geometry",
"error.gcv", "error.lepski")}, \code{ess.target = NULL} (NULL $\Rightarrow$ derive
from the global pick), \code{lepski.tau = 2}, \code{scale.grid = NULL} (NULL
$\Rightarrow$ the default of \S\ref{sec:grid}). \textbf{Every default reproduces
today's behavior bit-for-bit} (a single global $K$); the field mode is opt-in.
\item[I2.] \textbf{New resolver.} \code{.klp.resolve.prediction.support.size}
mirroring \code{.klp.resolve.prediction.chart.dim}: loops over anchors, runs
\S\ref{sec:algorithm}, returns \code{support.size.by.eval}. For the joint product,
a single \code{.klp.resolve.prediction.scale.and.dim} that returns both fields so the
$(K,d)$ choice is made together (one sort, one SVD sweep per anchor).
\item[I3.] \textbf{Prediction loop.} \code{.klp.predict.local.polynomial} already
loops per anchor and reads \code{chart.dim.by.eval[[i]]}; replace the scalar
\code{support.size} with \code{support.size.by.eval[[i]]}. Everything downstream
(neighborhood, weights, chart, fit) is already per-anchor.
\item[I4.] \textbf{Criterion helpers.} ESS from the kernel weights (reuse the Tier-1
ESS code) --- used both for the Mode~A target and for the \emph{ESS-based feasibility cap}
(\S\ref{sec:criteria}); Lepski from the $S_{i\cdot}$/SE machinery (the scale-field decision
rule); local GCV from the influence row in \code{R/lps\_uncertainty.R} as a
\emph{diagnostic only}.
\item[I5.] \textbf{Backend.} R backend only at first (as \code{local.auto} already
is). The native backends stay on the global path; document that the field mode forces
R.
\item[I6.] \textbf{Selection-aware bookkeeping (Mode B only).} When the scale is
error-driven, mark the fit as selection-dependent so the df/band code routes through
the selection-aware path; in Mode A, assert and record that $S$ stayed
response-independent. Additionally expose a \emph{generalized-df estimate} for Mode~B
--- the Ye perturbation / Efron covariance penalty, which \emph{reselects the scale and
refits} on each jittered $y$ (not a refit at frozen fields) --- so S3 and the band
reporting (D6) can report the selection cost $\mathrm{GDF}-\mathrm{tr}\,S(\hat K)$.
\end{enumerate}

%==============================================================================
\section{Synthetic test data}\label{sec:dgp}
%==============================================================================

Three generators give the answer keys the tests need. The first is new and must be
added to the frozen DGP library; the other two already exist.

\field{D-VARSCALE (new) --- a known scale change}
A 2-D surface in 3-D, sampled like the paraboloid generator (latent
$u\sim\text{Uniform}$ on the unit disk; $X=(u_1,u_2,(u_1^2+u_2^2)/(2R))$, curvature
knob $R$). The truth on the surface has a \emph{smooth half} and a \emph{wiggly half}:
\[
 f(u)=\sin(\pi\,g(u_1)\,u_1)\,\cos(\pi\,u_2),\qquad
 g(u_1)=1+\tfrac{k-1}{2}\big(1+\tanh(u_1/w)\big),
\]
so the oscillation frequency rises smoothly from $1$ on the left ($u_1\ll0$) to $k$ on
the right ($u_1\gg0$) across a transition of width $w$. Parameters (defaults):
$R=2$, $k=4$, $w=0.15$, $n=1500$, homoscedastic noise $\sigma=0.05$, recorded seed.
The \emph{known} property the tests use: the best local neighborhood is \emph{small}
on the right (fast oscillation) and \emph{large} on the left (smooth).

\smallskip\noindent\emph{No confound (the auditor checks this).} On the paraboloid, density
and curvature depend only on the \emph{radius} $\rho=\sqrt{u_1^2+u_2^2}$ (the surface area
element and the curvature are functions of $\rho$), hence are \emph{symmetric in $u_1$} and
identical on the two sides; only the frequency $g(u_1)$ is asymmetric in $u_1$. So the
smooth-vs-wiggly contrast is a \emph{pure frequency contrast}, not a density, boundary, or
curvature contrast. The generator records, per point: the side label, the radius $\rho$,
and a mask flagging the transition band ($|u_1|\le w$) and the domain boundary (an
outer-$\rho$ quantile), so the tracking gates (G3, G8) compare deep-smooth vs deep-wiggly
interior anchors only. The transition is excluded partly because the ramp's own derivative
adds local frequency there: the effective $u_1$-frequency of $f$ is
$\pi\big(g(u_1)+u_1\,g'(u_1)\big)$, and $g'$ peaks at the transition, so $|u_1|\le w$
carries extra wiggliness belonging to neither the deep-smooth nor the deep-wiggly regime. A
separate \emph{density-gradient variant} (more points where $u_1>0$, identical truth) is
provided for G8's positive control --- it is the only version that deliberately breaks the
density symmetry.

\field{G4 (existing) --- varying dimension}
The two-glued-strata generator: a $1$-D segment joined to a $2$-D patch, with the
true local dimension $1$ on one stratum and $2$ on the other, and region labels.

\field{G3a (existing) --- a clean curved surface}
The paraboloid with $f_{\text{smooth}}$, used as the homogeneous control (no real
scale or dimension change), to check LDS does \emph{not} invent variation.

%==============================================================================
\section{Test plan}\label{sec:tests}
%==============================================================================

Each test states its type, the data and configuration, the exact pass condition, and
the single change (``mutation'') that must make it fail. \textbf{GATE} $=$ a
deterministic pass/fail check; \textbf{STUDY} $=$ a decision rule on a measured
quantity. The implementer writes the tests; the auditor independently confirms each
mutation reddens the named test (the implementer does not run the mutations as
acceptance evidence).

\noindent\emph{What each mode is held to.} Mode~A (geometry-only) is tested for
\emph{reduction, linearity, well-posedness, and determinism}, and for adapting to
\emph{density/geometry but not function wiggliness} (G1, G2, G5, G6, G8). Mode~B--Lepski
is the mode tested for adapting to \emph{smooth-vs-wiggly truth} (G3, S1). Expecting
geometry-only to track frequency is a category error: on a DGP whose geometry is identical
on both sides, only an error-driven rule can see the frequency change.

\begin{gate}{G1 --- reduction to today (GATE)}
On G3a (homogeneous) with Mode A and $\mathrm{ESS}_{\text{target}}$ set to the global
pick's ESS, the chosen \code{support.size.by.eval} is (near) constant and its fitted
values match the current single-global-$K$ fit to algebraic tolerance $10^{-10}$ in
the degenerate case where the grid contains only the global $K$. \emph{Mutation:} make
the ESS target depend on the anchor index $\Rightarrow$ the field varies and the match
breaks.
\end{gate}

\begin{gate}{G2 --- linearity preserved in Mode A (GATE)}
In Mode A (geometry-only), the method is a linear smoother: extract $S$ by the
column-by-column probe (the E0.2 protocol) at the resolved fields, and confirm
$\hat y = S y$ to $10^{-10}$ for several random $y$, and $\mathrm{df}=\mathrm{tr}\,S$.
\emph{Mutation:} let the ESS target read $y$ (e.g.\ widen where $|y|$ is large)
$\Rightarrow$ $S$ becomes $y$-dependent and the identity fails. (This mutation is also
the proof that Mode~B is \emph{not} a fixed linear smoother --- expected, documented.)
\end{gate}

\begin{gate}{G3 --- the scale field tracks a known wiggliness change (Mode~B, GATE)}
On D-VARSCALE, under \textbf{Mode~B--Lepski}, the median chosen $K$ on the smooth (left)
side is materially larger than on the wiggly (right) side --- pin a margin, e.g.\ left
median $K \ge 1.5\times$ right median $K$ --- judged on interior anchors, excluding the
transition band ($|u_1|\le w$) and the domain boundary. \textbf{Mode~A is not expected to
pass this}, and must not: the geometry is identical on both sides (\S\ref{sec:dgp}), so
only an error-driven rule can see the frequency change --- the geometry-only behavior is
the control gate G8. \emph{Mutation:} replace the per-anchor Lepski choice with the global
$K$ $\Rightarrow$ the two sides' medians equalize and the check fails.
\end{gate}

\begin{gate}{G4 --- the dimension field tracks a known change (GATE)}
On the G4 generator, the per-anchor $d^\star$ equals $1$ on the $1$-D stratum and $2$ on
the $2$-D stratum for a large majority of interior anchors, \emph{excluding a transition
band around the glued-strata boundary} (where a neighborhood straddles both strata and the
local dimension is legitimately ambiguous); pin a fraction on the non-transition interior,
e.g.\ $\ge 0.9$, with the joint $(K,d)$ search. \emph{Mutation:} drop the eigengap branch
of the dimension rule $\Rightarrow$ the $1$-D stratum is over-estimated and the fraction
falls.
\end{gate}

\begin{gate}{G5 --- well-posedness always holds (GATE)}
At every anchor, $\binom{d^\star+p}{p}\le \mathrm{ESS}(K^\star)/\rho$ (the ESS-based
feasibility of \S\ref{sec:criteria}) and the local design's condition number $\le$
\code{ridge.condition.max}; on healthy data no anchor returns NA. \emph{Mutation:} switch
the feasibility test back to raw $K$ (or remove it) $\Rightarrow$ a heavily down-weighted
neighborhood selects a $(K,d)$ whose \emph{effective} points are too few and the design
becomes ill-conditioned.
\end{gate}

\begin{gate}{G6 --- determinism (GATE)}
Same data and seed $\Rightarrow$ bitwise-identical \code{support.size.by.eval} and
\code{chart.dim.by.eval}. \emph{Mutation:} introduce an unseeded tie-break in the
distance ordering $\Rightarrow$ two runs differ.
\end{gate}

\begin{gate}{G7 --- conditional linearity is exact even in Mode B (GATE)}
Run a Mode~B fit and \emph{freeze} the resolved fields $(K^\star_i,d^\star_i)$ it
produced. With those fields held fixed (no reselection), extract $S$ by the
column-by-column probe (the E0.2 protocol) and confirm $\hat y=S\,y$ to $10^{-10}$ for
several random $y$, with $\mathrm{df}=\mathrm{tr}\,S$. This isolates the claim of
\S\ref{sec:conditional}: the fit \emph{given the configuration} is exactly linear; the
only non-linearity lives in the selection map. \emph{Mutation:} re-resolve the fields
against each probe vector $y$ instead of freezing them $\Rightarrow$ the column probe no
longer reconstructs a single $S$ and the identity fails --- exhibiting precisely the
selection non-linearity that \S\ref{sec:conditional} describes.
\end{gate}

\begin{gate}{G8 --- geometry-only does not chase wiggliness (Mode~A control, GATE)}
On the same D-VARSCALE (uniform latent density, and curvature symmetric in $u_1$; only the
frequency $g(u_1)$ is asymmetric), \textbf{Mode~A}'s chosen $K$ is statistically flat across
the ramp: the smooth/wiggly side-median ratio lies in a tight band around $1$ (e.g.\
$[0.8,1.25]$) on interior anchors away from the boundary. This is the intended
\emph{negative} result --- a geometry-only rule cannot, and must not, read frequency; it is
the companion to G3.
\emph{Geometry-response check (in the right variable):} on a density-gradient variant of
D-VARSCALE (more points where $u_1>0$, identical truth), the adaptivity must appear in the
\emph{radius}, not the count. With a $K$-nearest-neighbor bandwidth ($h=$ the $K$-th-NN
distance), $\mathrm{ESS}(K)$ is nearly density-invariant, so the correct check is: $h_i$
shrinks in the denser region, $K_i$ stays $\approx$ stable, and $\mathrm{ESS}_i$ stays near
target. (Genuine $K$ movement, where it occurs, comes from the \emph{feasibility floor} ---
higher local dimension $\Rightarrow$ larger minimum $K$ --- which G4/G5 exercise; so the
flat-$K$ result here is the correct null, not an insensitive test.) \emph{Mutation:} let
Mode~A's ESS target read $y$ (widen where $|y|$ or the local oscillation is large)
$\Rightarrow$ the side ratio leaves the band --- the same response-leak G2 forbids, seen
here in the $K$-field.
\end{gate}

\begin{gate}{S1 --- accuracy benefit (STUDY)}
On D-VARSCALE (and as a secondary, G4), compare Truth-RMSE of LDS (Modes A and B)
against the global-$K$ baseline \emph{at matched df} (tune the baseline's $K$ and
LDS's target so both spend the same total df), paired on the same data and noise draw,
$R$ replicates. Pre-declared ``material'': median per-replicate
$\Delta\text{RMSE}>5\%$ in LDS's favor with a signed-rank $p<0.05$. Report both modes;
report where the gain concentrates (expected: the wiggly region and the dimension
boundary). \emph{Reporting allowance (LDS2):} matched df can be awkward to hit exactly for
adaptive fields, so the first LDS2 report may carry \emph{both} the raw selected-field fits
and the matched-df fits, with the \textbf{matched-df row as the decision} and the raw row as
context.
\end{gate}

\begin{gate}{S2 --- which mode, and what it does conditionally (STUDY)}
On the same data (LDS2), report Mode A vs Mode B: Truth-RMSE, the realized
\emph{conditional} df ($\mathrm{tr}\,S(\hat K)$), the selected scale/dimension summaries
(per-side $K$/$d$ medians and the $(K,d)$ fields), and --- for Mode B --- the honest
(training-removal / Lepski) versus shortcut LOO gap, so the bias of the cheap shortcut is
visible. The \emph{unconditional} cost (GDF versus $\mathrm{tr}\,S(\hat K)$) is
\emph{deferred to S3 / LDS3}, not reported here. This informs decision D1 (whether the
error-driven gain justifies leaving the exact-toolkit guarantee); the band-reporting choice
D6 is settled with S3 in LDS3.
\end{gate}

\begin{gate}{S3 --- selection-induced undercoverage, and its correction (STUDY)}
On a known-truth generator (D-VARSCALE, with G3a as the homogeneous control), over $R$
replicates with fresh noise, fit Mode~B and form pointwise bands two ways: the
\emph{conditional} band from $\mathrm{tr}\,S(\hat K)$ / $SS^\top$ at the selected
fields, and a \emph{selection-aware} band (the GDF / covariance-penalty-corrected df
from the I6 estimate, or a whole-pipeline bootstrap). Measure empirical coverage of the
true mean at the anchors. Pre-declared rule: the conditional band \emph{undercovers}
(coverage materially below nominal --- e.g.\ $\le 0.90$ for a nominal-$0.95$ band ---
concentrated in the wiggly region where reselection is most active), while the
selection-aware band restores coverage to within a pinned tolerance of nominal. Report
the realized $\mathrm{GDF}-\mathrm{tr}\,S(\hat K)$ gap. As a control, Mode~A's
conditional coverage should already sit near nominal (its only selection is the single
global target), quantifying that the two modes differ in exactly this respect.
\emph{Mutation (negative control):} feed the selection-aware path a frozen-field refit
(no reselection) in place of the reselect-and-refit estimate $\Rightarrow$ the corrected
df collapses back to $\mathrm{tr}\,S$ and the restored coverage disappears.
\emph{Phase \& cost:} this study is \emph{deferred to LDS3} (\S\ref{sec:sequence}) and is
\emph{not} an early blocking gate --- the GDF perturb-and-refit is computationally
expensive and statistically noisy, so it gates nothing in LDS0--LDS2; it informs D6 once
the cheaper checks pass.
\end{gate}

\begin{gate}{S4 --- oracle scale gap (STUDY, diagnostic)}
Using the \emph{truth} for evaluation only, compute the \emph{oracle} local scale map
$K^{\mathrm{orc}}_i=\arg\min_{K}\big(\hat m_K(c_i)-f(c_i)\big)^2$ (the size minimizing each
anchor's squared error against the known truth; average over noise draws to stabilize).
Report, per mode: the agreement between the chosen $K_i$ and $K^{\mathrm{orc}}_i$ (rank
correlation, and the smooth/wiggly side medians against the oracle's), and the Truth-RMSE
\emph{gap} to the oracle fit. This says how much of the available adaptivity each mode
recovers --- Mode~B--Lepski should track $K^{\mathrm{orc}}$; Mode~A should not (it has no
access to it), which is the quantitative form of G8. Diagnostic, not a blocking gate.
\end{gate}

\begin{gate}{S5 --- dimension-rule threshold sensitivity (STUDY, diagnostic)}
The dimension rule's $95\%$ energy and eigengap-$4$ thresholds (\S\ref{sec:dimrule}) are
reasonable first defaults, not derived constants. Sweep them --- energy
$\in\{0.90,0.95,0.99\}$, eigengap $\in\{2,4,8\}$ --- on the G4 generator (and D-VARSCALE)
and report the stability of the recovered dimension field and of Truth-RMSE. Soft
pass-for-robustness: the modal $d^\star$ per stratum is unchanged across the central part
of the sweep. This guards against the defaults being load-bearing flukes; sensitivity
diagnostic, not a blocking gate.
\end{gate}

\noindent\textbf{Reuse and conventions.} Tests use explicit fold ids and recorded
seeds; thresholds are pinned in the test files; the execution bundle follows the
Tier-0 pattern (clean committed tree, checksums, session info, realized quantities).
Mode~A tests assert the linear-smoother identity; Mode~B tests assert it does
\emph{not} hold and use the selection-aware error instead.

%==============================================================================
\section{Build-and-validation sequence (LDS0--LDS3)}\label{sec:sequence}
%==============================================================================

Implement and validate in four steps; each unlocks the next. This puts the cheap, decisive
checks first and the expensive selection-aware analysis last.

\begin{description}[leftmargin=0em,style=nextline]
\item[LDS0 --- geometry-only resolver.] Build the joint $(K_i,d_i)$ resolver in Mode~A and
prove \emph{reduction} (G1), \emph{linearity} (G2), \emph{determinism} (G6),
\emph{ESS-based well-posedness} (G5), and the \emph{geometry-not-wiggliness} control (G8).
No response-driven selection yet; this is the safe floor and the exact-toolkit baseline.
\item[LDS1 --- Mode~B--Lepski.] Add the Lepski scale rule and show it adapts to true local
smoothness on D-VARSCALE (G3), with conditional linearity at frozen fields still exact
(G7). Local GCV is wired only as a diagnostic.
\item[LDS2 --- comparison at matched df.] Compare global-$K$, Mode~A, and Mode~B at matched
\emph{conditional} df by Truth-RMSE (S1, S2), with the oracle gap (S4) and the
dimension-threshold sensitivity (S5) as diagnostics. \textbf{This is what decides the
shipped default (D1)} --- on accuracy at matched df, not on the unconditional cost.
\item[LDS3 --- selection-aware inference, and the smoothing question.] Only now study the
\emph{unconditional} cost (GDF / coverage, S3) and whether the raw $(K_i,d_i)$ fields need
spatial smoothing (the LDR / LDSR question, \S\ref{sec:program}). S3's perturb-and-refit is
expensive and noisy, so it lands here, never as an early gate.
\end{description}

%==============================================================================
\section{What the auditor checks}\label{sec:audit}
%==============================================================================

\begin{itemize}[leftmargin=1.4em]
\item Every mutation in \S\ref{sec:tests} reddens its named test (the central
deliverable); a test that stays green under its mutation is rejected.
\item Mode A genuinely never touches $y$ (read the criterion code, not just the
output): a gaussian fit's $S$ is response-independent (G2).
\item The honest LOO caveat is handled as one of the three stated mitigations, not
papered over (the local-GCV shortcut is not used as the headline error for $k$-NN).
\item The conditional-vs-unconditional distinction (\S\ref{sec:conditional}) is stated,
and any df/band claim says which it is: $\mathrm{tr}\,S(\hat K)$ is presented as the
\emph{conditional} df, never as the unconditional df of the pick-then-fit procedure.
G7 must show the fit is conditionally linear; S3 must show the conditional band
undercovers under Mode~B and the selection-aware correction restores coverage.
\item The GDF / covariance-penalty estimate (I6, S3) \emph{reselects} the scale on each
perturbed or bootstrapped $y$, not merely refits at frozen fields; an estimate that
freezes the fields measures $\mathrm{tr}\,S$, not GDF, and is rejected.
\item D-VARSCALE matches its specification (the frequency ramp $g(u_1)$, the curvature
$R$, the noise) \emph{and} the density/curvature symmetry in $u_1$ holds, so the ``known
scale change'' is a \emph{pure frequency contrast} (\S\ref{sec:dgp}); the G3/G8 tracking
gates exclude the transition band and boundary and are not vacuous.
\item The G3/G8 split is respected: Mode~A is judged by the geometry-only \emph{control}
(G8 --- flat $K$ across the frequency ramp), \emph{not} expected to track wiggliness; only
Mode~B--Lepski is held to G3. G8's positive control (density-gradient variant) confirms the
null is a real null, not an insensitive test.
\item Feasibility is on the \emph{effective} sample size,
$\mathrm{ESS}(K)\ge\rho\binom{d+p}{p}$ (not raw $K$); the G5 mutation (revert to raw $K$)
reddens. The ESS cap is \emph{gated to the LDS local-scale path}: the existing global-$K$
\code{chart.dim = "local.auto"} dimension is byte-for-byte unchanged (covered by the
bit-for-bit default pin).
\item Local GCV is treated as a \emph{diagnostic}, not the scale-field decision rule
(Lepski is); the build sequence (\S\ref{sec:sequence}) implements Lepski first.
\item The GDF / coverage study (S3) is \emph{not} used as an early blocking gate (deferred
to LDS3), and the dimension-rule defaults ($95\%$ energy, eigengap $4$) carry the S5
sensitivity diagnostic rather than being asserted as derived constants.
\item For G4, dimension recovery is judged \emph{off} the glued-strata transition band.
\item Defaults reproduce today's single-global-$K$ behavior bit-for-bit (the
backward-compatibility pin).
\item The accuracy study is paired and df-matched (not a comparison that lets the two
methods differ on flexibility as well as on locality).
\end{itemize}

%==============================================================================
\section{Decisions needed before implementation starts}\label{sec:decisions}
%==============================================================================

\begin{decide}
\textbf{D1.} Recommended local mode and shipped default: the API default stays today's
global $K$ bit-for-bit (I1); when per-point scale is enabled, the recommended front-runner
is Mode~B--Lepski (local error-driven, stabilized) with Mode~A as the exact-toolkit
fallback/floor, and the \emph{shipped} default is decided empirically by S1/S2 at matched
df (accuracy gain) in phase LDS2, with the GDF/coverage cost (S3) a later LDS3 check.
Confirm, or keep Mode~A as the default.\\
\textbf{D2.} $\mathrm{ESS}_{\text{target}}$ for Mode A: derive from the global pick
[recommended, gives the clean reduction], or pin a fixed value?\\
\textbf{D3.} The candidate scale grid (the \S\ref{sec:grid} default, or a pinned set).\\
\textbf{D4.} Mode~B--GCV honest-LOO mitigation (GCV is now a \emph{diagnostic},
\S\ref{sec:criteria}; Lepski is the decision rule): true leave-one-out on scored points,
leave-a-block-out, or Lepski-only [Lepski recommended for the scale field].\\
\textbf{D5.} Sequencing with the next phase: does the regularization phase
\emph{absorb} scale (one joint dimension-and-scale field, LDSR, built directly on
LDS), or land dimension-only regularization (LDR) first and add scale (LDSR) after?\\
\textbf{D6.} Band/df reporting under selection (\S\ref{sec:conditional}): report
\emph{conditional} bands on the selected scale, clearly labeled [recommended default],
or report \emph{unconditional} (selection-aware) bands via the GDF / covariance-penalty
correction or an outer split? Mode~A may stay conditional (its only selection is the
single global target); Mode~B should at least \emph{report} the GDF gap (S3) even if the
default band stays conditional.\\
\textbf{D7.} ESS feasibility factor $\rho$ in $\mathrm{ESS}(K)\ge\rho\binom{d+p}{p}$
(\S\ref{sec:criteria}): default $\rho=1.5$ [recommended --- effective points comfortably
above the parameter count], $\rho=1$ (bare feasibility), or another pinned value.\\
\textbf{D8.} Lepski noise estimate $\hat\sigma$ and tolerance $\tau$: global robust (pilot
global-$K$ residual scale) [recommended first default], local, CV-derived, or
known-in-synthetic, with $\tau=2$ default. Lepski's selected scale is sensitive to both, so
both must be pinned before LDS1.
\end{decide}

%==============================================================================
\section{How LDS fits the wider program}\label{sec:program}
%==============================================================================

LDS produces the two raw per-anchor fields $(K_i,d_i)$. The next phases smooth them
over the chart nerve: \textbf{LDR} (dimension only, scale held fixed) and
\textbf{LDSR} (both fields, co-regularized). LDS is therefore the prerequisite for
both, and is independent of the prediction-synchronization axis (PS), so it can be
built in parallel with PS once the LPS base (Tier 0--4) closes.

Two sibling plans are to be written next and can proceed in parallel with LDS, since
they share the chart machinery but address different objects: (i) \textbf{local
covariance (LCov)} --- the second-order analogue, which will reuse this same
dimension/scale ladder; and (ii) the \textbf{16S stratification-deconvolution} data
models --- the realistic synthetic generators and the component/stratum recovery that
the capstone consumes. This document is scoped to LDS only; it is written so those two
can reference its ladder rather than re-derive it.
\end{document}
